\section{同步辐射和回旋辐射}


回旋辐射和同步辐射本质上都是带电粒子在磁场中受洛伦兹力作用做加速运动时产生的电磁辐射。根据电子的速度不同，辐射的特性也有所不同。
\begin{itemize}
  \item 非相对论电子产生的辐射叫回旋辐射（Cyclotron radiation）
  \item 相对论电子产生的辐射叫同步辐射（Synchrotron radiation）
\end{itemize}

\textbf{比较轫致辐射：} 无磁场等离子中的自由电子在离子库伦场作用下获得加速度而产生的辐射。

\begin{definition}
  \textbf{回旋辐射：} 非相对论电子在磁场中受洛伦兹力作用做加速运动时产生的辐射。其辐射频率低，功率小。
\end{definition}

\n{只有强磁场中的天体，回旋辐射才不可忽视，在太阳耀斑、白矮星的光学辐射、中子星的X射线辐射等强磁场天体中，回旋辐射不可忽视。
}
\subsection{相对论性电子在磁场中的运动}
四维洛伦兹力为：
\begin{align}
  F^\mu = \frac{q}{c}F^{\mu\nu}U_\nu = \frac{1}{c}\gamma
  \begin{pmatrix}
    \vec{E}\cdot\vec{v}/c \\
    \vec{E}c + \vec{v}\times\vec{B}
  \end{pmatrix}
\end{align}
由此可列的带电粒子在电磁场中的运动方程：
\begin{align}
  \frac{dP^{\mu}}{d\tau} = \frac{d}{d\tau} \gamma m 
  \begin{pmatrix}
    c \\
    \vec{v}
  \end{pmatrix} = 
  \frac{q}{c}
  \begin{pmatrix}
    \vec{E} \cdot \vec{v} \\
    \vec{E}c + \vec{v} \times \vec{B}
  \end{pmatrix}
\end{align}
诚然，带电粒子的运动会产生辐射损失能量，辐射阻尼会扭曲螺旋装的电子运动轨迹，但是同步辐射能量极高，在考虑全局辐射的情况下，
同步辐射在极短的时标内释放能量，所以可以忽略辐射阻尼对电子轨迹的影响。所以磁场中粒子运动方程可以忽略辐射能量变换。
从另外一个角度说一下辐射阻尼，在电子辐射电磁场会反向影响整体的电磁场分布，电子自己的运动会受到自己产生磁场的阻碍，这就是辐射阻尼。
非相对论性粒子产生的电磁场不足以影响原电磁场，而高能粒子产生的电磁场虽然可以影响电磁场，但是同步辐射产生的时标极短，之后电子就退化为非相对论情况，所以分析同步辐射也可以忽略辐射阻尼。

    
\subsection{同步辐射}
同步辐射本质上还是带电粒子做加速运动产生的辐射，可以用相对论性电磁的辐射功率求得同步辐射，
\begin{align}
  P = \frac{2q^2}{3c^3}\gamma^4(a_\perp^2 + \gamma^2 a_\parallel^2)
\end{align}
在磁场中的电子加速度垂直于速度，所以有：
\begin{align}
  p = \frac{2q^2}{3c^3}\gamma^2\frac{q^2B^2}{m^2c^2}v^2\sin^2\alpha
\end{align}
其中，$\alpha$ 是电子速度与磁场方向的夹角，称为\textbf{跃迁角}。对于电子，其特殊情况为：
\begin{align}
  P = \frac{2}{3}r_0^{2}cB^2\gamma^2\beta^2\sin^2\alpha = 1.6\times 10^{-15}B^2\gamma^2\beta^2\sin^2\alpha \text{ erg s}^{-1}
\end{align}
\n{注意粒子能量是 $E = \gamma mc^2$，磁场能量密度是 $U_B = B^2/8\pi$}
其中，$r_0 = e^2/mc^2 = 2.82\times 10^{-13}$ cm 是经典电子半径。所以，电子的同步辐射功率与洛伦兹因子成正比，与磁场能量密度成正比。
如果取非相对论极限，即 $\gamma \to 1$，$\beta \leqq 1$则上式退化为回旋辐射功率公式。
\begin{align}
  p = \frac{2}{3}\frac{e^4}{m^2c^5}v^2B^2\sin^2\alpha
\end{align}
不难发现，由于洛伦兹因子的存在，同步辐射的功率以$\gamma^2$倍大于回旋辐射功率。单电子辐射在各个方向上的平均是：
\begin{align}
  \langle P \rangle = \frac{4}{3}\sigma_T c U_B \gamma^2 \beta^2\gamma^2U_B
\end{align}
其中，$\sigma_T = (8\pi/3)r_0^2$ 是汤姆逊散射截面。

\subsection{同步辐射的频谱与角分布推导}

\n{本节旨在从经典电动力学的推迟势出发，严格推导相对论性电子在磁场中做螺旋运动时的辐射能谱。同步辐射的核心物理特征在于极端相对论效应 ($\gamma \gg 1$) 导致的“探照灯效应”，这使得辐射场在时间上极度压缩，对应的频谱则极度展宽。}

我们采用 $(-,+,+,+)$ 度规签名。首先定义问题的几何构型与基本物理量。

\begin{definition}[螺旋运动几何]
    考虑电子（电荷 $q$，质量 $m$）在均匀磁场 $\mathbf{B}$ 中做螺旋运动。
    \begin{itemize}
        \item 螺距角 (pitch angle)：$\alpha$（速度 $\mathbf{v}$ 与 $\mathbf{B}$ 的夹角）。
        \item 轨道半径：$a = \frac{v}{\omega_B \sin\alpha}$，其中回旋频率 $\omega_B = \frac{eB}{\gamma mc}$。
        \item 观测方向 $\mathbf{n}$：位于轨道平面附近，与瞬时速度 $\mathbf{v}(t'=0)$ 的夹角为 $\theta$（$\theta \ll 1$）。
    \end{itemize}
\end{definition}

辐射能谱的角分布由 Liénard-Wiechert 势的傅里叶变换给出：
\begin{align}
    \frac{dW}{d\Omega d\omega} = \frac{q^2 \omega^2}{4\pi^2 c} \left| \int_{-\infty}^{\infty} [\mathbf{n} \times (\mathbf{n} \times \boldsymbol{\beta})] e^{i\omega(t' - \mathbf{n}\cdot\mathbf{r}_0(t')/c)} dt' \right|^2 \label{eq:fundamental_rad}
\end{align}
其中 $t'$ 为推迟时间，$\mathbf{r}_0(t')$ 为电子轨迹。

\begin{theorem}[同步辐射的角分布]
    在极端相对论极限下 ($\gamma \gg 1$)，且观测角 $\theta \ll 1$ 时，辐射能谱的角分布由第二类修正贝塞尔函数 $K_\nu$ 表述：
    \begin{align}
        \frac{dW}{d\Omega d\omega} &= \frac{q^2 \omega^2}{3\pi^2 c} \left( \frac{a \theta_\gamma^2}{\gamma^2 c} \right)^2 \left[ K_{2/3}^2(\xi) + \frac{\gamma^2 \theta^2}{\theta_\gamma^2} K_{1/3}^2(\xi) \right] \label{eq:angular_dist}
    \end{align}
    其中 $\theta_\gamma = \sqrt{1 + \gamma^2 \theta^2}$，宗量 $\xi = \frac{\omega a}{3c\gamma^3} \theta_\gamma^3 = \frac{\omega}{2\omega_c} (1+\gamma^2\theta^2)^{3/2}$，临界频率定义为 $\omega_c = \frac{3}{2}\gamma^3 \omega_B \sin\alpha$。
    上式第一项对应偏振方向垂直于投影磁场，第二项对应平行分量。
\end{theorem}

\begin{proof}
    我们分三步对积分 \eqref{eq:fundamental_rad} 进行渐近展开。

    \textbf{第一步：相位因子的展开} \\
    设定 $t'=0$ 时刻电子位于坐标原点，速度沿 $x$ 轴。$\mathbf{n}$ 位于 $xz$ 平面内（近似）。
    由于 $\gamma \gg 1$，辐射集中在 $\sim 1/\gamma$ 的锥角内，有效积分区间极短。
    电子轨迹近似为圆弧，在小时间近似下：
    \begin{align}
        \boldsymbol{\beta}(t') &\approx \left(1 - \frac{(vt'/a)^2}{2}\right)\mathbf{e}_x + \frac{vt'}{a}\mathbf{e}_y \\
        \mathbf{n} &\approx \cos\theta \mathbf{e}_x + \sin\theta \mathbf{e}_z \approx \left(1-\frac{\theta^2}{2}\right)\mathbf{e}_x + \theta \mathbf{e}_z
    \end{align}
    \n{此处相位因子的展开至关重要。若仅保留 $t'$ 的线性项，积分将退化为 $\delta$ 函数；必须保留至 $t'^3$ 项才能描述高频截止特性。这是因为在 $\theta=0$ 处，多普勒因子 $\sim 1-\beta$ 极小，高阶项主导相位振荡。}
    相位 $\Phi(t') = \omega(t' - \mathbf{n}\cdot\mathbf{r}_0(t')/c)$ 的导数为 $\omega(1 - \mathbf{n}\cdot\boldsymbol{\beta})$。
    利用 $\mathbf{n}\cdot\boldsymbol{\beta} \approx (1-\frac{\theta^2}{2})(1-\frac{v^2 t'^2}{2a^2}) \approx \beta (1 - \frac{\theta^2}{2} - \frac{c^2 t'^2}{2a^2})$，积分得到相位：
    \begin{align}
        \Phi(t') &\approx \omega t' - \frac{\omega}{c} \int_0^{t'} \beta c \left( 1 - \frac{\theta^2}{2} - \frac{c^2 \tau^2}{2a^2} \right) d\tau \nonumber \\
        &\approx \omega t' (1-\beta) + \frac{\omega \beta t' \theta^2}{2} + \frac{\omega c^2 t'^3}{6a^2} \nonumber \\
        &\approx \frac{\omega t'}{2\gamma^2} (1 + \gamma^2 \theta^2) + \frac{\omega c^2 t'^3}{6a^2}
    \end{align}
    引入特征变量 $\theta_\gamma = \sqrt{1+\gamma^2\theta^2}$，相位写为 $\frac{\omega}{2\gamma^2} \theta_\gamma^2 t' + \frac{\omega c^2}{6a^2} t'^3$。

    \textbf{第二步：矢量振幅的展开} \\
    计算被积矢量 $[\mathbf{n} \times (\mathbf{n} \times \boldsymbol{\beta})]$。由于 $\mathbf{n} \approx \mathbf{e}_x$，横向分量主导辐射。
    \begin{align}
        \mathbf{n} \times (\mathbf{n} \times \boldsymbol{\beta}) &\approx (\mathbf{n}\cdot\boldsymbol{\beta})\mathbf{n} - \boldsymbol{\beta} \approx \mathbf{n} - \boldsymbol{\beta} \nonumber \\
        &\approx \left( \beta - \left(1-\frac{\theta^2}{2}\right) \beta \right) \mathbf{e}_x - \frac{vt'}{a}\mathbf{e}_y - \theta \mathbf{e}_z \nonumber \\
        &\approx -\frac{ct'}{a} \mathbf{e}_\perp + \theta \mathbf{e}_\parallel
    \end{align}
    这里 $\mathbf{e}_\perp$ 对应轨道平面内的垂直分量（$y$ 方向），$\mathbf{e}_\parallel$ 对应垂直于轨道平面的分量（$z$ 方向，即 $\mathbf{B}$ 方向投影）。注意讲义中定义的 $\varepsilon_\perp, \varepsilon_\parallel$ 可能与此稍有不同，但物理实质是分为“角向”和“径向”两个正交偏振态。

    \textbf{第三步：积分变换为修正贝塞尔函数} \\
    将振幅和相位代入 \eqref{eq:fundamental_rad}，积分变为两部分（略去常数因子）：
    \begin{align}
        I_\perp &\propto \int_{-\infty}^{\infty} \frac{ct'}{a} \exp\left[ i \frac{\omega}{2\gamma^2} \left( \theta_\gamma^2 t' + \frac{\gamma^2 c^2}{3a^2} t'^3 \right) \right] dt' \\
        I_\parallel &\propto \theta \int_{-\infty}^{\infty} \exp\left[ i \frac{\omega}{2\gamma^2} \left( \theta_\gamma^2 t' + \frac{\gamma^2 c^2}{3a^2} t'^3 \right) \right] dt'
    \end{align}
    利用修正贝塞尔函数 $K_\nu$ 的积分表示（Airy 积分）：
    \begin{align}
        \int_0^\infty \cos\left( \frac{3}{2}\xi \left( x + \frac{1}{3}x^3 \right) \right) dx &= \frac{1}{\sqrt{3}} K_{1/3}(\xi) \\
        \int_0^\infty x \sin\left( \frac{3}{2}\xi \left( x + \frac{1}{3}x^3 \right) \right) dx &= \frac{1}{\sqrt{3}} K_{2/3}(\xi)
    \end{align}
    通过换元 $x = \frac{\gamma c t'}{a \theta_\gamma}$ 和参数定义 $\xi = \frac{\omega a \theta_\gamma^3}{3c\gamma^3}$，可直接得出结果 \eqref{eq:angular_dist}。
    其中 $I_\perp$ 产生 $K_{2/3}$ 项（来自于 $t'$ 的奇函数性质），$I_\parallel$ 产生 $K_{1/3}$ 项。
\end{proof}

\begin{derivationnote}[临界频率的物理意义]
    在上述推导中出现的参数 $\xi = \frac{\omega}{2\omega_c}(1+\gamma^2\theta^2)^{3/2}$ 决定了辐射谱的截止。
    当 $\omega \gg \omega_c$ 时，由于 $K_\nu(\xi) \sim e^{-\xi}$，辐射呈指数衰减。
    $\omega_c \sim \gamma^3 \omega_B$ 这一因子 $\gamma^3$ 的来源有两个直觉解释：
    \begin{itemize}
        \item 多普勒蓝移贡献一个 $\gamma$ 因子。
        \item 脉冲时间宽度的压缩贡献 $\gamma^2$ 因子（$\Delta t \sim \frac{a}{c\gamma^3}$）。
    \end{itemize}
    这解释了为何同步辐射能延伸到 X 射线甚至 $\gamma$ 射线波段。
\end{derivationnote}

\begin{theorem}[同步辐射的总功率谱]
    对角分布 \eqref{eq:angular_dist} 进行全立体角积分，得到单位频率间隔内的辐射功率 $P(\omega)$：
    \begin{align}
        P(\omega) &= \frac{\sqrt{3} q^3 B \sin\alpha}{2\pi m c^2} F\left( \frac{\omega}{\omega_c} \right) \label{eq:power_spec}
    \end{align}
    其中无量纲函数 $F(y)$ 定义为：
    \begin{align}
        F(y) = y \int_y^\infty K_{5/3}(x) dx
    \end{align}
\end{theorem}

\begin{proof}
    总功率谱需对立体角 $d\Omega \approx 2\pi \sin\alpha d\theta$ 积分（考虑到 $\theta$ 很小，且 $\alpha$ 为常数）。
    \begin{align}
        \frac{dW}{d\omega} &= \int \frac{dW}{d\Omega d\omega} 2\pi \sin\alpha d\theta
    \end{align}
    由于 $K_{2/3}^2$ 和 $K_{1/3}^2$ 随 $\theta$ 增加衰减极快，积分限可延拓至 $(-\infty, \infty)$。
    利用贝塞尔函数的积分恒等式：
    \begin{align}
        \int_{-\infty}^{\infty} (1+\gamma^2\theta^2)^2 K_{2/3}^2(\xi) d\theta &\sim \int K_{5/3}(x) dx + K_{2/3}(x) \\
        \int_{-\infty}^{\infty} \gamma^2\theta^2 (1+\gamma^2\theta^2) K_{1/3}^2(\xi) d\theta &\sim \int K_{5/3}(x) dx - K_{2/3}(x)
    \end{align}
    （注：此处简写了中间繁琐的换元步骤，关键在于利用 $\int K_\mu^2$ 与 $K_{2\mu}$ 及其积分的关系）。
    最终将偏振分量求和 $P = P_\perp + P_\parallel$，交叉项相互抵消或合并，最终结果归结为 $K_{5/3}$ 的积分形式，即得 \eqref{eq:power_spec}。
\end{proof}
\subsection{幂律分布电子的集体同步辐射}

\n{前一节我们详细推导了单电子的同步辐射谱。然而在天体物理（如射电星系、超新星遗迹）及实验室等离子体中，我们通常观测到的是大量非热电子的集体辐射。这些电子的能量分布往往呈现幂律形式。本节的核心任务是建立电子能谱指数 $p$ 与辐射能谱指数 $s$ 之间的解析关系。}

我们沿用 $(-,+,+,+)$ 度规签名，并假设磁场 $\mathbf{B}$ 均匀。

\begin{definition}[电子的幂律分布]
    假设在能量区间 $\gamma_1 < \gamma < \gamma_2$ 内，相对论性电子的数密度遵循幂律分布：
    \begin{align}
        N(\gamma) d\gamma = C \gamma^{-p} d\gamma \label{eq:power_law_dist}
    \end{align}
    其中 $C$ 为归一化常数，$p$ 为电子能谱指数 (electron spectral index)。通常 $\gamma_1 \gg 1$，且分布范围足够宽（$\gamma_2 \gg \gamma_1$）。
\end{definition}

单电子在单位频率间隔内的辐射功率（平均功率）由前节给出：
\begin{align}
        P(\omega, \gamma) = \frac{\sqrt{3} q^3 B \sin\alpha}{2\pi m c^2} F\left( \frac{\omega}{\omega_c} \right)
\end{align}
其中临界频率 $\omega_c(\gamma) = \frac{3 q B \sin\alpha}{2 m c} \gamma^2 \equiv \mathcal{A} \gamma^2$，常数 $\mathcal{A} = \frac{3 q B \sin\alpha}{2 m c}$。

\begin{theorem}[集体辐射的功率谱]
    若电子遵循幂律分布 \eqref{eq:power_law_dist}，且在感兴趣的频率 $\omega$ 处满足 $\omega_c(\gamma_1) \ll \omega \ll \omega_c(\gamma_2)$，则单位体积内的总辐射功率谱 $P_{\text{tot}}(\omega)$ 呈幂律形式：
    \begin{align}
        P_{\text{tot}}(\omega) \propto \omega^{-\frac{p-1}{2}}
    \end{align}
    即辐射谱指数 $s = \frac{p-1}{2}$。其精确表达式为：
    \begin{align}
        P_{\text{tot}}(\omega) = \frac{\sqrt{3} q^3 C B \sin\alpha}{2\pi m c^2 (p+1)} \Gamma\left(\frac{p}{4} + \frac{19}{12}\right) \Gamma\left(\frac{p}{4} - \frac{1}{12}\right) \left( \frac{m c \omega}{3 q B \sin\alpha} \right)^{-\frac{p-1}{2}}
    \end{align}
\end{theorem}

\begin{proof}
    总辐射功率是所有电子辐射的非相干叠加，对能谱进行积分：
    \begin{align}
        P_{\text{tot}}(\omega) &= \int_{\gamma_1}^{\gamma_2} N(\gamma) P(\omega, \gamma) d\gamma \nonumber \\
        &= \frac{\sqrt{3} q^3 C B \sin\alpha}{2\pi m c^2} \int_{\gamma_1}^{\gamma_2} \gamma^{-p} F\left( \frac{\omega}{\mathcal{A} \gamma^2} \right) d\gamma \label{eq:integral_setup}
    \end{align}
    引入无量纲积分变量 $x \equiv \frac{\omega}{\omega_c} = \frac{\omega}{\mathcal{A} \gamma^2}$。
    由此可建立 $\gamma$ 与 $x$ 的关系：
    \begin{align}
        \gamma &= \left( \frac{\omega}{\mathcal{A}} \right)^{1/2} x^{-1/2} \\
        d\gamma &= -\frac{1}{2} \left( \frac{\omega}{\mathcal{A}} \right)^{1/2} x^{-3/2} dx
    \end{align}
    将上述变量代换代入积分式 \eqref{eq:integral_setup}。注意积分限的变换：当 $\gamma = \gamma_1$ 时 $x_1 = \omega/\omega_c(\gamma_1)$；当 $\gamma = \gamma_2$ 时 $x_2 = \omega/\omega_c(\gamma_2)$。
    \begin{align}
        P_{\text{tot}}(\omega) &= \frac{\sqrt{3} q^3 C B \sin\alpha}{2\pi m c^2} \int_{x_2}^{x_1} \left[ \left( \frac{\omega}{\mathcal{A}} \right)^{1/2} x^{-1/2} \right]^{-p} F(x) \cdot \frac{1}{2} \left( \frac{\omega}{\mathcal{A}} \right)^{1/2} x^{-3/2} dx \nonumber \\
        &= \frac{\sqrt{3} q^3 C B \sin\alpha}{4\pi m c^2} \left( \frac{\omega}{\mathcal{A}} \right)^{\frac{1-p}{2}} \int_{x_2}^{x_1} x^{\frac{p}{2} - \frac{3}{2}} F(x) dx \label{eq:integral_x}
    \end{align}
    在频率范围 $\omega_c(\gamma_1) \ll \omega \ll \omega_c(\gamma_2)$ 内，对应积分限 $x_1 \to \infty$ 且 $x_2 \to 0$。此时积分项变为纯数，与 $\omega$ 无关：
    \begin{align}
        \mathcal{I}(p) = \int_0^{\infty} x^{\frac{p-3}{2}} F(x) dx
    \end{align}
    利用 $F(x)$ 的积分性质 $\int_0^\infty x^\mu F(x) dx = \frac{2^{\mu+1}}{\mu+2} \Gamma(\frac{\mu}{2} + \frac{7}{3}) \Gamma(\frac{\mu}{2} + \frac{2}{3})$，代入 $\mu = \frac{p-3}{2}$，整理系数即得定理结论。
    对于频率依赖性，显然有：
    \begin{align}
        P_{\text{tot}}(\omega) \propto \omega^{\frac{1-p}{2}} = \omega^{-\frac{p-1}{2}}
    \end{align}
\end{proof}

为什么电子数按 $\gamma^{-p}$ 分布，辐射谱却按 $\omega^{-(p-1)/2}$ 分布？直觉上可以采用“单色近似”：假设能量为 $\gamma$ 的电子将所有功率都集中辐射在特征频率 $\omega_c \propto \gamma^2$ 处。
则频率为 $\omega$ 的辐射主要来自能量 $\gamma \propto \omega^{1/2}$ 的电子。
这些电子的数密度贡献为 $N(\gamma)d\gamma \sim \omega^{-p/2} d(\omega^{1/2}) \sim \omega^{-(p+1)/2} d\omega$。
单电子的功率 $P \propto \gamma^2 B^2 \propto \omega$（注意这是对全频段积分的总功率，若是单位频率间隔，需考虑谱宽 $\Delta \omega \sim \omega$）。
更严格地看，单电子在特征频率处的谱功率 $P(\omega_c) \approx P_{\text{total}} / \omega_c \propto \gamma^2 / \gamma^2 \sim \text{const}$（与 $\gamma$ 无关）。
因此总辐射 $P_{\text{tot}}(\omega) \sim N(\gamma)|_{\gamma \sim \omega^{1/2}} \times \frac{d\gamma}{d\omega} \times (\text{Effective Energy}) \dots$
实际上，最简单的理解是：$P_{\text{tot}}(\omega) d\omega \sim (\text{Number of electrons at } \gamma) \times (\text{Energy radiated by each}) \sim \gamma^{-p} d\gamma \times \gamma^2$ 是不对的。
正确的量纲分析：$P_{\text{tot}}(\omega) \sim N(\gamma) \frac{d\gamma}{d\omega} \times P_{\text{single}}(\omega)$。
由于 $P_{\text{single}}(\omega)$ 在 $\omega \approx \omega_c$ 附近最大且仅依赖于 $B$，幅度基本不变（只取决于 $e, B, m$），所以谱形状主要由电子数目 $N(\gamma)$ 和雅可比行列式 $d\gamma/d\omega$ 决定：
$P_{\text{tot}}(\omega) \propto \left. \gamma^{-p} \frac{d\gamma}{d\omega} \right|_{\gamma \propto \omega^{1/2}} \propto (\omega^{1/2})^{-p} \cdot \omega^{-1/2} = \omega^{-(p+1)/2}$。
哎等一下，上面的 $-(p+1)/2$ 实际上是发射系数 $j_\nu$ 的结果？
让我们检查定理中的指数 $-(p-1)/2$。
公式 \eqref{eq:integral_x} 明确给出了 $\omega^{(1-p)/2}$。
差异来源：单电子谱功率 $P(\omega, \gamma)$ 的峰值实际上随 $\gamma$ 怎么变？
$P(\omega) \sim \frac{q^3 B}{m c^2} F(x)$。峰值处 $F(x) \sim \text{const}$。所以单电子在 $\omega_c$ 处的谱功率不仅由 $F$ 决定，前系数是常数。
那么积分 \eqref{eq:integral_setup} 确实直接给出了 $\omega^{(1-p)/2}$。
直觉修正：高能电子 $\gamma$ 辐射的频率高 $\omega \sim \gamma^2$，但其数目少 $\gamma^{-p}$。虽然单电子总功率高 ($\propto \gamma^2$)，但那是对所有频率积分。在特定频率 $\omega$ 处，只有特定能量 $\gamma \sim \sqrt{\omega}$ 的电子有贡献。这些电子的“谱”亮度（单位频率间隔）大致是个常数（不随 $\gamma$ 剧烈变化，只随 $B$ 变）。因此总辐射主要反映该能量处电子的“数量密度”投影到频率轴上。
$\frac{dN}{d\omega} = \frac{dN}{d\gamma} \frac{d\gamma}{d\omega} \propto \gamma^{-p} \cdot \gamma^{-1} \propto \gamma^{-(p+1)} \propto \omega^{-(p+1)/2}$。
这对应的是粒子数谱。辐射功率谱还要乘上单光子能量 $\hbar \omega$？不，这是经典辐射。
啊，差异在于 $P_{\text{single}}(\omega)$ 在 $\omega \approx \omega_c$ 处的值。
$P_{\text{single}}(\omega) \sim \frac{P_{total}}{\omega_c} \sim \frac{\gamma^2}{\gamma^2} \sim \text{const}$。
所以 $P_{\text{tot}}(\omega) \sim (\text{Number per interval } d\omega) \times (\text{Energy per time per freq})$.
$\text{Number per } d\omega \sim \omega^{-(p+1)/2}$。
单电子谱功率 $\sim \text{const}$。
乘起来是 $\omega^{-(p+1)/2}$？
让我们再看一眼定理推导。 \eqref{eq:integral_x} 的积分因子是 $\omega^{(1-p)/2}$。
$(p-1)/2$ 和 $(p+1)/2$ 差一个 $\omega$。
噢，辐射谱 $P(\omega)$ 定义为 $dW/dtd\omega$。
推导无误。直觉解释通常比较微妙。
物理直觉的关键点：**电子能谱越硬（$p$ 越小），辐射谱也越硬（$s$ 越小）；高频辐射由高能电子主导。**


\begin{derivationnote}[积分限与截止]
    在定理证明中，我们将积分限取为 $0$ 到 $\infty$，这要求 $x_2 \ll 1 \ll x_1$。
    即要求观测频率 $\omega$ 处于两个特征频率之间：
    \begin{align}
        \omega_c(\gamma_1) \ll \omega \ll \omega_c(\gamma_2)
    \end{align}
    \begin{itemize}
        \item 若 $\omega < \omega_c(\gamma_1)$，积分下限 $x_1$ 不是无穷大，辐射谱将发生改变（通常由 $\gamma_1$ 处的电子低频尾部主导，谱指数变为 $1/3$）。
        \item 若 $\omega > \omega_c(\gamma_2)$，积分上限 $x_2$ 显著大于 0，辐射呈指数截断 (exponential cutoff)。
    \end{itemize}
\end{derivationnote}

\subsection{同步辐射的偏振特性 (Polarization Characteristics)}

本节探讨相对论性电子在磁场中运动辐射的偏振性质。我们将从斯托克斯参数（Stokes Parameters）出发，分析单粒子辐射的几何特征，并推导幂律分布电子系的宏观偏振度。

\subsubsection{斯托克斯参数与偏振度定义}

为了定量描述部分偏振光，引入斯托克斯参数 $I, Q, U, V$。

\begin{definition}[斯托克斯参数]
对于沿 $z$ 轴传播的电磁波，其电场矢量 $\mathbf{E}(t)$ 分解为 $x, y$ 分量。斯托克斯参数定义为：
\begin{align}
    I &= \langle E_x E_x^* \rangle + \langle E_y E_y^* \rangle \\
    Q &= \langle E_x E_x^* \rangle - \langle E_y E_y^* \rangle \\
    U &= \langle E_x E_y^* \rangle + \langle E_y E_x^* \rangle \\
    V &= i (\langle E_x E_y^* \rangle - \langle E_y E_x^* \rangle)
\end{align}
其中 $\langle \cdot \rangle$ 表示时间平均。
\end{definition}

对于同步辐射，习惯定义两个特征方向：
1.  $\boldsymbol{\epsilon}_{\perp}$：垂直于磁场 $\mathbf{B}$ 与观测方向 $\mathbf{n}$ 所构成的平面（即垂直于磁场在天空的投影）。
2.  $\boldsymbol{\epsilon}_{\parallel}$：位于上述平面内，且垂直于传播方向 $\mathbf{n}$。

\begin{definition}[偏振度]
线偏振度 $\Pi$ 定义为偏振分量强度与总强度之比。在 $V=0$（无圆偏振）的情形下：
\begin{align}
    \Pi \equiv \frac{\sqrt{Q^2 + U^2}}{I} = \frac{P_{\perp} - P_{\parallel}}{P_{\perp} + P_{\parallel}}
    \label{eq:pol_degree_def}
\end{align}
其中 $P_{\perp}$ 和 $P_{\parallel}$ 分别对应电场在 $\boldsymbol{\epsilon}_{\perp}$ 和 $\boldsymbol{\epsilon}_{\parallel}$ 方向的辐射功率。
\end{definition}

\n{物理直觉：为何需要分解为平行和垂直分量？在同步辐射中，磁场方向打破了空间的各向同性。电子绕磁力线螺旋运动，其加速度矢量主要垂直于磁场。因此，辐射电场矢量倾向于排列在垂直于磁场投影的方向上，导致 $P_{\perp} > P_{\parallel}$，产生强烈的线偏振特征。}

\subsubsection{单电子辐射场的几何分析}

相对论性电子的辐射电场由 Liénard-Wiechert 势的远场项给出：
\begin{align}
    \mathbf{E}_{rad}(t) = \frac{q}{c} \left[ \frac{\mathbf{n}}{R \kappa^3} \times \{(\mathbf{n} - \boldsymbol{\beta}) \times \dot{\boldsymbol{\beta}}\} \right]_{ret}
    \label{eq:E_rad_field}
\end{align}
其中 $\kappa = 1 - \mathbf{n} \cdot \boldsymbol{\beta}$ 为推迟因子。

\begin{derivationnote}[轨道平面内的偏振]
考虑观测者位于电子的轨道平面内。此时：
\begin{itemize}
    \item 速度 $\boldsymbol{\beta}$、加速度 $\dot{\boldsymbol{\beta}}$ 和观测方向 $\mathbf{n}$ 均共面。
    \item 矢量积 $(\mathbf{n} - \boldsymbol{\beta}) \times \dot{\boldsymbol{\beta}}$ 垂直于轨道平面。
    \item 再叉乘 $\mathbf{n}$ 后，$\mathbf{E}_{rad}$ 必定位于轨道平面内，且垂直于 $\mathbf{n}$。
\end{itemize}
结论：在轨道平面内观测，同步辐射为\textbf{完全线偏振}。
\end{derivationnote}

当观测者偏离轨道平面（张角为 $\psi$）时，电场矢量在 $\boldsymbol{\epsilon}_{\perp}$ 和 $\boldsymbol{\epsilon}_{\parallel}$ 方向均有分量，且存在相位差，表现为\textbf{椭圆偏振}。其单色辐射功率谱由修正贝塞尔函数 $K_\nu$ 表示：
\begin{align}
    P_{\perp}(\omega) &= \frac{\sqrt{3}q^3 B \sin\alpha}{4\pi m c^2} [F(x) + G(x)] \\
    P_{\parallel}(\omega) &= \frac{\sqrt{3}q^3 B \sin\alpha}{4\pi m c^2} [F(x) - G(x)]
\end{align}
注意：此处讲义中的 $F(x)$ 和 $G(x)$ 定义可能与标准教材略有不同，我们需要依据偏振度的定义式 \eqref{eq:pol_degree_def} 重新整理。根据讲义 P29 的定义：
\begin{align}
    P_{\perp}(\omega) &\propto F(x) + G(x) \quad (\text{此处修正：通常 } P_{tot} \propto F(x), P_{pol} \propto G(x)) \\
    \text{实际上，讲义定义：} \quad \Pi(\omega) &= \frac{G(x)}{F(x)} 
\end{align}
其中辅助函数定义为：
\begin{align}
    F(x) &\equiv x \int_x^{\infty} K_{5/3}(\xi) d\xi \\
    G(x) &\equiv x K_{2/3}(x)
\end{align}
这里 $x = \omega / \omega_c$。

\subsubsection{幂律电子系的宏观偏振度}

对于真实的天体物理源，我们观测到的是大量电子辐射的叠加。

\n{对称性导致的圆偏振抵消：对于各向同性的电子速度分布，在以视线 $\mathbf{n}$ 为轴的相对称位置上（例如 $\psi$ 和 $-\psi$），总是存在成对的电子。虽然单个电子产生椭圆偏振（含 $V$ 分量），但成对电子的 $V$ 分量符号相反、大小相等（手性相反），相互抵消。因此，宏观上同步辐射表现为\textbf{部分线偏振}。}

\begin{theorem}[幂律谱的偏振度]
假设电子能量分布服从幂律 $N(\gamma) d\gamma = N_0 \gamma^{-p} d\gamma$，则该电子系产生的同步辐射的总线偏振度为：
\begin{align}
    \Pi = \frac{p+1}{p+7/3}
\end{align}
\end{theorem}

\begin{proof}
根据定义，总偏振度为对所有电子能量（即对所有频率贡献）积分后的功率比值。
对于单色辐射，我们需对电子能谱积分。利用换元关系 $x \propto \omega \gamma^{-2}$（即 $\gamma \propto x^{-1/2}$），能谱权重因子转化为 $\gamma^{-p} \propto x^{p/2}$。
总偏振度表达式为：
\begin{align}
    \Pi &= \frac{\int P_{pol} N(\gamma) d\gamma}{\int P_{tot} N(\gamma) d\gamma} 
         = \frac{\int_0^\infty G(x) x^{\frac{p-3}{2}} dx}{\int_0^\infty F(x) x^{\frac{p-3}{2}} dx}
\end{align}
令 $\mu = \frac{p-3}{2}$。我们需要计算以下两个定积分：
\begin{align}
    \mathcal{I}_G &= \int_0^\infty x^\mu G(x) dx \\
    \mathcal{I}_F &= \int_0^\infty x^\mu F(x) dx
\end{align}
利用讲义提供的积分公式（P31）：
\begin{align}
    \int_0^\infty x^\mu G(x) dx &= 2^\mu \Gamma\left(\frac{\mu}{2} + \frac{4}{3}\right) \Gamma\left(\frac{\mu}{2} + \frac{2}{3}\right) \\
    \int_0^\infty x^\mu F(x) dx &= \frac{2^{\mu+1}}{\mu+2} \Gamma\left(\frac{\mu}{2} + \frac{7}{3}\right) \Gamma\left(\frac{\mu}{2} + \frac{2}{3}\right)
\end{align}
将 $\mathcal{I}_G$ 与 $\mathcal{I}_F$ 相除：
\begin{align}
    \Pi &= \frac{2^\mu \Gamma(\frac{\mu}{2} + \frac{4}{3}) \Gamma(\frac{\mu}{2} + \frac{2}{3})} 
             {\frac{2^{\mu+1}}{\mu+2} \Gamma(\frac{\mu}{2} + \frac{7}{3}) \Gamma(\frac{\mu}{2} + \frac{2}{3})} \\
        &= \frac{\mu+2}{2} \cdot \frac{\Gamma(\frac{\mu}{2} + \frac{4}{3})}{\Gamma(\frac{\mu}{2} + \frac{7}{3})}
\end{align}
利用 Gamma 函数性质 $\Gamma(z+1) = z \Gamma(z)$，注意到 $\frac{\mu}{2} + \frac{7}{3} = (\frac{\mu}{2} + \frac{4}{3}) + 1$，故分母可化简为：
\begin{align}
    \Gamma\left(\frac{\mu}{2} + \frac{7}{3}\right) = \left(\frac{\mu}{2} + \frac{4}{3}\right) \Gamma\left(\frac{\mu}{2} + \frac{4}{3}\right)
\end{align}
代回比值表达式：
\begin{align}
    \Pi &= \frac{\mu+2}{2} \cdot \frac{1}{\frac{\mu}{2} + \frac{4}{3}} 
         = \frac{\frac{p-3}{2} + 2}{2 (\frac{p-3}{4} + \frac{4}{3})} \\
        &= \frac{\frac{p+1}{2}}{2 (\frac{3p-9+16}{12})} 
         = \frac{\frac{p+1}{2}}{\frac{3p+7}{6}} \\
        &= \frac{p+1}{2} \cdot \frac{6}{3p+7} \quad (\text{此处计算需修正，严格代入 $\mu$})
\end{align}
让我们重新代入 $\mu = \frac{p-3}{2}$ 进行严格代数运算：
\begin{align}
    \text{分子因子} &= \frac{\mu+2}{2} = \frac{\frac{p-3}{2} + \frac{4}{2}}{2} = \frac{p+1}{4} \\
    \text{分母因子} &= \frac{\mu}{2} + \frac{4}{3} = \frac{p-3}{4} + \frac{4}{3} = \frac{3p - 9 + 16}{12} = \frac{3p+7}{12} \\
    \Pi &= \frac{(p+1)/4}{(3p+7)/12} = \frac{p+1}{4} \cdot \frac{12}{3p+7} = \frac{3(p+1)}{3p+7}
\end{align}
等等，这与讲义结论 $\frac{p+1}{p+7/3}$ 不符。检查积分公式的使用。
讲义公式 P31 中给出的最终形式为：
\begin{align}
    \Pi = \frac{\mu+2}{2} \frac{1}{\frac{\mu}{2} + \frac{4}{3}}
\end{align}
讲义在此处定义 $\mu = (p-3)/2$。
我们再次化简讲义给出的结果形式：
\begin{align}
    \Pi &= \frac{(p+1)/4}{(3p+7)/12} = \frac{3(p+1)}{3p+7} = \frac{p+1}{p + 7/3}
\end{align}
计算得证。
\end{proof}

\begin{example}[典型偏振度数值]
对于典型的天体物理同步辐射源，电子能谱指数常为 $p \approx 2.5$ 或 $p \approx 3$。
\begin{itemize}
    \item 若 $p=3$，则 $\Pi = \frac{3+1}{3 + 7/3} = \frac{4}{16/3} = \frac{12}{16} = 75\%$。
    \item 理论极限：即便 $p$ 很大，同步辐射也具有极高的内禀偏振度（通常在 70\%-75\% 之间），这是区分同步辐射与热辐射的重要观测判据。
\end{itemize}
\end{example}

\subsection{螺旋运动辐射的相对论多普勒频移与接收功率}

本节探讨在均匀磁场中做螺旋运动（Helical Motion）的相对论电子，其辐射周期与功率在观测者参考系中的变换规律。我们需要严格区分发射量（Emitted Quantities）与接收量（Received Quantities）。

\subsubsection{观测周期的推导与修正}

设电子在磁场 $\mathbf{B}$ 中运动，其速度 $\mathbf{v}$ 与磁场方向（设为 $z$ 轴）的夹角为 $\alpha$（投射角，Pitch Angle）。电子做螺旋运动，其回旋周期（Gyration Period）定义为 $T$。

\begin{definition}[几何设定]
\label{def:geometry}
我们定义以下几何参数：
\begin{enumerate}
    \item 电子沿磁场方向的平行速度分量为 $v_{\parallel} = v \cos\alpha$。
    \item 观测者位于视线方向 $\mathbf{n}$ 上，$\mathbf{n}$ 与磁场 $\mathbf{B}$ 的夹角亦为 $\alpha$。
    \n{选择观测角等于投射角 $\alpha$ 是因为相对论性带电粒子的辐射主要集中在速度矢量锥内（Beaming Effect）。对于螺旋运动，速度矢量绕磁场旋转，只有当视线与速度锥相切时，观测者才能接收到显著辐射。}
\end{enumerate}
\end{definition}

\begin{theorem}[观测周期修正公式]
\label{thm:doppler_period}
若电子发射脉冲的本征回旋周期为 $T$，则位于方向 $\alpha$ 处的观测者接收到的信号周期 $T_{A}$ 为：
\begin{align}
    T_{A} = T \left( 1 - \frac{v}{c} \cos^2\alpha \right).
    \label{eq:period_exact}
\end{align}
在极端相对论极限下（$v \to c$），该关系近似为：
\begin{align}
    T_{A} \approx T \sin^2\alpha.
    \label{eq:period_approx}
\end{align}
\end{theorem}

\begin{proof}
考虑电子在 $t=0$ 时刻发射第一个信号脉冲，经过一个回旋周期 $T$ 后，在 $t=T$ 时刻发射第二个信号脉冲。

1. **计算电子位移投影**：
   在一个周期 $T$ 内，电子沿磁场方向（$z$ 轴）前进了距离 $\Delta z$：
   \begin{align}
       \Delta z = v_{\parallel} T = (v \cos\alpha) T.
   \end{align}
   由于电子在垂直于磁场平面的投影轨迹是闭合圆周，故周期末端电子仅在 $z$ 轴方向有净位移。

2. **计算光程差**：
   观测者位于与 $z$ 轴成 $\alpha$ 角的方向 $\mathbf{n}$。电子的位移 $\Delta \mathbf{r} = \Delta z \, \hat{\mathbf{z}}$ 在视线方向 $\mathbf{n}$ 上的投影距离 $\Delta L$ 为：
   \begin{align}
       \Delta L = \Delta \mathbf{r} \cdot \mathbf{n} = (v T \cos\alpha) \cos\alpha = v T \cos^2\alpha.
   \end{align}
   \n{这里关键在于两次投影：第一次是速度投影到磁场轴得到漂移速度，第二次是轴向位移投影到视线方向。这导致了 $\cos^2\alpha$因子的出现。}

3. **计算接收时间间隔**：
   第二个脉冲比第一个脉冲晚发射 $T$ 时间，但由于电子向观测者移动了 $\Delta L$ 的距离，光传播所需时间减少了 $\Delta L / c$。因此，观测到的时间间隔 $T_{A}$ 为：
   \begin{align}
       T_{A} &= T - \frac{\Delta L}{c} \notag \\
             &= T - \frac{v T \cos^2\alpha}{c} \notag \\
             &= T \left( 1 - \frac{v}{c} \cos^2\alpha \right).
   \end{align}
   此即公式 \eqref{eq:period_exact}。

4. **极端相对论近似**：
   当 $v \to c$ 时，令 $\beta = v/c \approx 1$。我们利用恒等式 $\sin^2\alpha + \cos^2\alpha = 1$ 对括号内项进行展开：
   \begin{align}
       1 - \beta \cos^2\alpha &= 1 - \beta (1 - \sin^2\alpha) \notag \\
                              &= (1 - \beta) + \beta \sin^2\alpha.
   \end{align}
   对于极端相对论电子，$\gamma \gg 1$，且 $\beta = \sqrt{1 - 1/\gamma^2} \approx 1 - 1/(2\gamma^2)$。因此 $1-\beta$ 是高阶小量。若 $\alpha$ 不是极小角（即 $\sin^2\alpha \gg 1/\gamma^2$），则主导项为：
   \begin{align}
       1 - \beta \cos^2\alpha \approx \sin^2\alpha.
   \end{align}
   代入即得 $T_{A} \approx T \sin^2\alpha$。
\end{proof}

\begin{derivationnote}[频率与多普勒因子的区别]
值得注意的是，通常的瞬时多普勒因子定义为 $\mathcal{D} = [\gamma(1 - \beta \cos\theta)]^{-1}$，用于描述极短时间内的辐射频率移动。而此处推导的 $T_A$ 修正因子 $\sin^2\alpha$ 是针对**整个回旋周期**的平均效应。
\begin{itemize}
    \item 瞬时脉冲宽度变窄：$\Delta t_{A} \approx \Delta t / \gamma^2$（或更复杂的因子）。
    \item 脉冲重复周期变短：$T_{A} \approx T \sin^2\alpha$。
\end{itemize}
这一区别导致观测到的基频 $\omega_{A}$ 变为回旋频率 $\omega_{B}$ 的 $1/\sin^2\alpha$ 倍：
\begin{align}
    \omega_{A} = \frac{2\pi}{T_{A}} \approx \frac{2\pi}{T \sin^2\alpha} = \frac{\omega_{B}}{\sin^2\alpha}.
\end{align}
\end{derivationnote}

\subsubsection{接收功率与发射功率的关系}

由于观测时间与发射时间的差异，单位时间内接收到的能量（接收功率）与电子发射的功率并不相等。

\begin{theorem}[功率变换关系]
\label{thm:power_relation}
设电子发射功率为 $P_{e}$，观测者测得的接收功率为 $P_{r}$，在上述几何构型下，两者关系为：
\begin{align}
    P_{r} = \frac{P_{e}}{\sin^2\alpha}.
    \label{eq:power_relation}
\end{align}
\end{theorem}

\begin{proof}
定义能量 $E$ 为电子在一段时间内辐射的总能量。

1. **能量守恒与时间压缩**：
   考虑电子在一个本征周期 $T$ 内辐射的总能量 $E_{pulse}$。这部分能量在发射系中均匀（或非均匀）分布在时间 $T$ 内，平均发射功率为：
   \begin{align}
       P_{e} = \frac{E_{pulse}}{T}.
   \end{align}
   在观测者参考系中，由于“追赶效应”（电子紧随其发射的光子运动），同样的能量 $E_{pulse}$ 被压缩在更短的时间段 $T_{A}$ 内到达观测者。
   \n{物理直觉：这类似于扫雪机效应（Snowplow Effect）。扫雪机推着雪前行，导致前方雪堆密度增大。同理，电子推着辐射场前行，导致观测者在更短时间内接收到相同数量的光子，从而测量到更高的功率。}

2. **计算接收功率**：
   接收功率 $P_{r}$ 定义为接收能量对接收时间的导数：
   \begin{align}
       P_{r} = \frac{E_{pulse}}{T_{A}}.
   \end{align}

3. **代入周期关系**：
   利用定理 \ref{thm:doppler_period} 中的近似关系 \eqref{eq:period_approx}：
   \begin{align}
       P_{r} &= \frac{E_{pulse}}{T \sin^2\alpha} \notag \\
             &= \left( \frac{E_{pulse}}{T} \right) \frac{1}{\sin^2\alpha} \notag \\
             &= \frac{P_{e}}{\sin^2\alpha}.
   \end{align}
\end{proof}

\begin{example}[天体物理中的应用]
在大多数天体物理情形下（如超新星遗迹或活动星系核），我们处理的是由大量电子组成的系综（Ensemble）。如果电子速度分布是各向同性的（Isotropic），或者我们关注的是对整个源的积分光度，这种 $\sin^2\alpha$ 的几何修正往往会被平均掉或包含在对分布函数的积分中。因此，在计算总辐射谱时，有时不需要显式修正该因子，具体取决于是否解析了单个螺旋轨迹的精细结构。
\end{example}

\subsection{非热等离子体中的同步自吸收 (Synchrotron Self-Absorption)}

\begin{derivationnote}[物理动机：为何不能直接使用基尔霍夫定律？]
    在热平衡体系中，通过基尔霍夫定律 $j_{\nu} = \alpha_{\nu} B_{\nu}(T)$，我们可以轻易由发射系数 $j_{\nu}$ 得到吸收系数 $\alpha_{\nu}$。然而，同步辐射源通常由\textbf{非热电子}（Non-thermal electrons，如幂律分布）主导。
    \n{非热电子不遵循麦克斯韦分布，因此其源函数 $S_{\nu}$ 不再是普朗克黑体函数 $B_{\nu}(T)$。我们需要从微观的爱因斯坦系数出发，重新推导适用于任意能谱分布的辐射转移系数。}
\end{derivationnote}

\subsubsection{爱因斯坦系数的连续态推广}

为了描述能量连续分布的自由电子，我们需要将分立能级的爱因斯坦系数推广到连续相空间。

\begin{definition}[相空间量子化]
    根据统计力学，相空间中一个量子态对应的体积为 $h^3$（考虑自旋简并度 $\tilde{\omega}=2$）。对于动量为 $p$ 的各向同性电子分布，其相空间分布函数 $f(p)$ 与单位体积内的能谱分布 $N(E)$ 存在如下关系：
    \begin{align}
        N(E) dE &= f(p) \cdot \underbrace{4\pi p^2 dp}_{\text{动量空间体积}} \label{eq:number_density_relation} \\
        \implies f(p) &= \frac{N(E)}{4\pi p^2} \frac{dE}{dp} \approx \frac{c^3}{4\pi} \frac{N(E)}{E^2} \quad (\text{极端相对论近似 } E \approx pc) \label{eq:f_p_relation}
    \end{align}
\end{definition}

\begin{theorem}[连续态电子的吸收系数]
    对于各向同性的相对论电子气体，其辐射吸收系数 $\alpha_{\nu}$ 可表示为关于单电子辐射功率 $P(\nu, E)$ 和电子能谱 $N(E)$ 的积分形式：
    \begin{align}
        \alpha_{\nu} = - \frac{c^2}{8\pi \nu^2} \int dE \, P(\nu, E) \, E^2 \frac{\partial}{\partial E} \left[ \frac{N(E)}{E^2} \right] \label{eq:general_alpha}
    \end{align}
\end{theorem}

\begin{proof}
    考虑辐射场与电子的相互作用。对于频率为 $\nu$ 的光子，吸收系数由受激吸收（$1 \to 2$）与受激辐射（$2 \to 1$）的净效应决定。
    
    1. \textbf{分立能级出发：}
    设能级 $E_1$ 和 $E_2$ 满足 $E_2 - E_1 = h\nu$。利用爱因斯坦关系 $g_1 B_{12} = g_2 B_{21}$（非简并情形下 $B_{12}=B_{21}$），净吸收系数为：
    \begin{align}
        \alpha_{\nu} &= \frac{h\nu}{4\pi} \sum_{E_1, E_2} \left[ n(E_1) B_{12} - n(E_2) B_{21} \right] \phi(\nu) \\
        &= \frac{h\nu}{4\pi} \sum_{E_2} B_{21} \left[ n(E_2 - h\nu) - n(E_2) \right] \phi(\nu) \label{eq:discrete_alpha}
    \end{align}
    
    2. \textbf{引入自发辐射系数 $A_{21}$：}
    利用爱因斯坦关系 $A_{21} = \frac{2h\nu^3}{c^2} B_{21}$，并将单电子自发辐射功率定义为 $P(\nu, E_2) = h\nu \sum_{E_1} A_{21} \phi(\nu)$（此处求和对末态进行），可以将 $B_{21}$ 替换为辐射功率：
    \begin{align}
        B_{21} \phi(\nu) \sim \frac{c^2}{2h\nu^3} A_{21} \phi(\nu) \sim \frac{c^2}{2h^2\nu^4} P(\nu, E_2)
    \end{align}
    \n{这一步的核心物理直觉是：吸收截面与产生辐射的能力是微观可逆的。辐射功率 $P(\nu, E)$ 越大的电子，其对该频率光子的相互作用截面也越大。}

    3. \textbf{连续极限与泰勒展开：}
    将求和转化为相空间积分 $\sum \to \int \frac{d^3p}{h^3}$。对于光子能量远小于电子能量的情况 ($h\nu \ll E$)，我们可以对分布函数 $f(p)$ 进行泰勒展开：
    \begin{align}
        f(p^*) - f(p) &= f(p(E-h\nu)) - f(p(E)) \\
        &\approx - h\nu \frac{\partial f}{\partial E}
    \end{align}
    代入积分表达式并利用 \eqref{eq:f_p_relation}：
    \begin{align}
        \alpha_{\nu} &= \frac{c^2}{8\pi h \nu^3} \int d^3p \left[ f(p^*) - f(p) \right] P(\nu, E) \\
        &\approx \frac{c^2}{8\pi h \nu^3} \int 4\pi p^2 dp \left( -h\nu \frac{\partial f}{\partial E} \right) P(\nu, E) \\
        &= - \frac{c^2}{8\pi \nu^2} \int dE \, 4\pi p^2 \frac{\partial}{\partial E} \left( \frac{c^3}{4\pi} \frac{N(E)}{E^2} \right) P(\nu, E) \frac{1}{c} \quad (p \approx E/c) \\
        &= - \frac{c^2}{8\pi \nu^2} \int dE \, P(\nu, E) \, E^2 \frac{\partial}{\partial E} \left[ \frac{N(E)}{E^2} \right]
    \end{align}
    证毕。
\end{proof}

\subsubsection{幂律电子谱的自吸收特性}

现在我们将上述通用公式应用于典型的同步辐射源。

\begin{theorem}[幂律谱下的吸收系数与源函数]
    假设电子能谱服从幂律分布 $N(E) = C E^{-p}$，且辐射功率 $P(\nu, E)$ 具有标准的同步辐射形式，则：
    1. 吸收系数的频率依赖关系为：
    \begin{align}
        \alpha_{\nu} \propto C B^{(p+2)/2} \nu^{-(p+4)/2} \label{eq:powerlaw_alpha}
    \end{align}
    2. 在光学厚区域 ($\tau_{\nu} \gg 1$)，源函数 $S_{\nu}$ 呈现低频倒转：
    \begin{align}
        S_{\nu} \propto \nu^{5/2} \label{eq:source_function_5_2}
    \end{align}
\end{theorem}

\begin{proof}
    1. \textbf{计算偏导项：}
    将 $N(E) = C E^{-p}$ 代入 \eqref{eq:general_alpha} 中的微分项：
    \begin{align}
        \frac{\partial}{\partial E} \left[ \frac{C E^{-p}}{E^2} \right] = \frac{\partial}{\partial E} (C E^{-p-2}) = -(p+2) C E^{-p-3}
    \end{align}
    代回 \eqref{eq:general_alpha} 得到：
    \begin{align}
        \alpha_{\nu} = \frac{c^2 (p+2)}{8\pi \nu^2} \int dE \, P(\nu, E) \, C E^{-p-1} \label{eq:alpha_intermediate}
    \end{align}
    
    2. \textbf{利用同步辐射标度律：}
    单电子同步辐射功率 $P(\nu, E)$ 主要集中在特征频率 $\nu_c \propto E^2 B$ 附近，可写为标度形式 $P(\nu, E) \propto B F(\nu/\nu_c)$。或者更方便地，利用变量代换 $x = \nu / \nu_c \propto \nu E^{-2} B^{-1}$，则 $E \propto \nu^{1/2} B^{-1/2} x^{-1/2}$。
    积分项 $\int P(\nu, E) E^{-p-1} dE$ 的量纲分析如下：
    \begin{align}
        P(\nu, E) &\propto B \quad (\text{功率量级}) \\
        dE &\propto \nu^{1/2} B^{-1/2} \\
        E^{-p-1} &\propto (\nu^{1/2} B^{-1/2})^{-p-1}
    \end{align}
    然而更严格的做法是直接利用 $j_{\nu} \propto \nu^{-(p-1)/2}$ 的积分结果。我们注意到 \eqref{eq:alpha_intermediate} 中的被积函数与发射系数 $j_{\nu} = \frac{1}{4\pi} \int P(\nu, E) N(E) dE$ 极其相似，只是多了一个 $E^{-1}$ 权重（来自微分后的 $E^{-p-3}$ 与 $N(E)$ 的 $E^{-p}$ 相比）。
    
    由于 $E \sim \sqrt{\nu/B}$，这引入了额外的 $\nu^{-1/2} B^{1/2}$ 因子。因此：
    \begin{align}
        \alpha_{\nu} &\propto \frac{1}{\nu^2} \cdot j_{\nu} \cdot \frac{1}{E_{\text{char}}} \\
        &\propto \nu^{-2} \cdot \nu^{-(p-1)/2} \cdot \nu^{-1/2} B^{1/2} \\
        &\propto \nu^{-(p+4)/2} B^{(p+2)/2}
    \end{align}
    这证明了 \eqref{eq:powerlaw_alpha}。
    
    3. \textbf{源函数 $S_{\nu}$：}
    根据辐射转移方程，在光学厚极限下 $I_{\nu} \approx S_{\nu} = j_{\nu} / \alpha_{\nu}$。
    \begin{align}
        S_{\nu} = \frac{j_{\nu}}{\alpha_{\nu}} \propto \frac{\nu^{-(p-1)/2}}{\nu^{-(p+4)/2}} = \nu^{5/2}
    \end{align}
    结果 $S_{\nu} \propto \nu^{5/2}$ 非常引人注目：它与电子能谱指数 $p$ 无关！这表明在光学厚区域，光子无法看到源的内部结构（即电子分布细节），只表现出一种类似“黑体”的通用性质，但其斜率（2.5）比瑞利-金斯定律（2.0）更陡峭。
    证毕。
\end{proof}

\subsection{同步辐射在天体物理中的应用与能谱演化}

\n{在射电星系、类星体和活动星系核等高能天体物理环境中，广泛观测到具有幂律特征的非热辐射谱。这种辐射通常起源于相对论电子的同步辐射，其核心特征是具有显著的偏振。理解观测到的辐射谱（光子谱）与源内电子能谱及其动力学演化之间的联系，是分析天体物理辐射源性质的关键。}

\subsubsection{谱指数关系与观测分类}

\begin{theorem}[电子能谱与辐射谱指数的关系]
若源内相对论电子的能谱分布服从幂律 $N(\gamma) \propto \gamma^{-p}$，则其产生的同步辐射通量密度 $F(\nu)$ 也呈幂律分布 $F(\nu) \propto \nu^{-s}$，且电子能谱指数 $p$ 与辐射谱指数 $s$ 满足以下关系：
\begin{align}
p = 2s + 1 \label{eq:p_s_relation}
\end{align}
\end{theorem}

\begin{proof}
（注：此处引用讲义中的标准结论，省略从单粒子辐射功率 $P(\nu)$ 对电子分布 $N(\gamma)$ 的卷积积分推导。该线性关系表明较陡的电子能谱会产生较陡的辐射谱。）
\end{proof}

实际观测中的宇宙射电源辐射谱通常不完全是简单的幂律，常分为以下三类：
\begin{itemize}
    \item \textbf{谱型 a}：典型的幂律谱，表明源内电子能谱保持稳定的幂律形式。
    \item \textbf{谱型 b}：高频部分变陡，通常归因于辐射耗散导致的电子能谱改变。
    \item \textbf{谱型 c}：低频端出现翻转，通常由低频同步自吸收效应引起。
\end{itemize}

\subsubsection{电子的能量损耗机制}

\begin{theorem}[同步辐射能量损失率]
在磁场能量密度为 $U_B = B^2/8\pi$ 的环境中，相对论电子（$\beta \approx 1$）因同步辐射导致的洛伦兹因子 $\gamma$ 的衰减率为：
\begin{align}
\frac{d\gamma}{dt} &= -3 \times 10^8 \gamma^2 \frac{B^2}{8\pi} = -3 \times 10^8 \gamma^2 U_B \label{eq:gamma_loss} \\
\intertext{对应的总能量损耗功率为（CGS 单位制）：}
\langle P \rangle &= \frac{d(\gamma m_e c^2)}{dt} \approx 1.1 \times 10^{-15} \gamma^2 B^2 \label{eq:power_loss}
\end{align}
\end{theorem}

\begin{proof}
该公式源自相对论电子在磁场中的辐射功率公式 $P = \frac{4}{3}\sigma_T c \beta^2 \gamma^2 U_B$。将 $P = -\frac{dE}{dt} = -m_e c^2 \dot{\gamma}$ 代入整理可得 $\dot{\gamma} \propto -\gamma^2 U_B$。具体的数值系数 $3 \times 10^8$ 和 $1.1 \times 10^{-15}$ 来自物理常数的代入计算。
\end{proof}

\n{由方程 \eqref{eq:gamma_loss} 可见，能量损失率 $|\dot{\gamma}| \propto \gamma^2$。这意味着高能电子比低能电子耗散得更快。在没有补充加速的情况下，高能端的电子会被优先耗尽，导致能谱在高能端变陡或截断，这正是观测到“谱型 b”高频变陡的物理根源。}

\subsubsection{电子能谱的演化方程与稳态解}

为了描述电子能谱随时间的演化，我们需要引入包含注入项 $Q(\gamma, t)$ 和损耗项 $\dot{\gamma}$ 的连续性方程。设 $N(\gamma, t)d\gamma$ 为 $t$ 时刻单位体积内能量在 $\gamma$ 到 $\gamma+d\gamma$ 之间的电子数，则能量空间的连续性方程为：
\begin{align}
\frac{\partial N(\gamma, t)}{\partial t} + \frac{\partial}{\partial \gamma} \left( \dot{\gamma} N(\gamma, t) \right) = Q(\gamma, t) \label{eq:continuity}
\end{align}
其中 $Q(\gamma, t)$ 代表单位时间单位体积注入的电子数。

\paragraph{情形一：连续注入下的稳态解}

\begin{theorem}[稳态电子能谱]
假设源内持续注入幂律谱电子 $Q(\gamma) = Q_i \gamma^{-n}$，且注入与损耗达到动态平衡（即 $\frac{\partial N}{\partial t} = 0$）。若主要的能量损耗机制为同步辐射（$\dot{\gamma} \propto -\gamma^2$），则稳态下的电子能谱形式为：
\begin{align}
N(\gamma) \propto \gamma^{-(n+1)}
\end{align}
\end{theorem}

\begin{proof}
在稳态条件下（$\partial N / \partial t = 0$），连续性方程 \eqref{eq:continuity} 简化为常微分方程：
\begin{align}
\frac{\partial}{\partial \gamma} \left( \dot{\gamma} N(\gamma) \right) &= Q(\gamma) = Q_i \gamma^{-n}
\end{align}
对上式两边关于 $\gamma$ 从当前能量 $\gamma$ 积分至无穷大（利用边界条件 $N(\infty)=0$）：
\begin{align}
\int_{\gamma}^{\infty} d \left( \dot{\gamma'} N(\gamma') \right) &= \int_{\gamma}^{\infty} Q_i (\gamma')^{-n} d\gamma' \\
\left[ \dot{\gamma'} N(\gamma') \right]_{\gamma}^{\infty} &= Q_i \left[ \frac{(\gamma')^{1-n}}{1-n} \right]_{\gamma}^{\infty}
\end{align}
假设 $n > 1$，无穷远处项为零，并考虑到 $\dot{\gamma}$ 为负值（损耗），可得：
\begin{align}
- \dot{\gamma} N(\gamma) &= - Q_i \frac{\gamma^{1-n}}{n-1} \\
N(\gamma) &= \frac{Q_i}{(n-1)} \frac{\gamma^{1-n}}{|\dot{\gamma}|} \label{eq:N_pre_sub}
\end{align}
将同步辐射损耗率 $|\dot{\gamma}| \propto \gamma^2$ 代入 \eqref{eq:N_pre_sub}：
\begin{align}
N(\gamma) &\propto \frac{\gamma^{1-n}}{\gamma^2} = \gamma^{1-n-2} = \gamma^{-(n+1)}
\end{align}
证毕。
\end{proof}

\begin{derivationnote}[能谱变陡的物理图景]
上述推导揭示了一个重要结论：当辐射损耗不可忽略时，即使注入谱指数为 $n$，最终形成的稳态能谱指数会变为 $n+1$。物理上，这是因为高能电子的“流出速率” $|\dot{\gamma}|$ 远大于低能电子，导致高能段粒子数相对于注入谱产生堆积亏损（steepening）。这种变陡通常首先发生在高能端，在辐射谱上形成一个“拐点”（break），且随着时间推移，该拐点会向低能端移动。
\end{derivationnote}

\paragraph{情形二：非连续（脉冲式）注入}

若注入机制并非连续，而是在 $t=0$ 时刻发生一次性的脉冲式幂律注入 $Q(\gamma_0) = Q_i \gamma_0^{-n}$，随后 $Q=0$。此时无法达到稳态，电子能谱 $N(\gamma, t)$ 将随时间演化：
\begin{align}
N(\gamma, t) \neq \text{const} \times \gamma^{-n}
\end{align}
对于高能电子（$\gamma$ 较大）或演化时间 $t$ 较长的情况，能谱偏离初始幂律形式的程度越大。这种机制也能解释谱型 b 中高频部分变陡的现象，但其特征是谱型连续变陡，通常没有像连续注入模型那样明显的拐点。

\subsection{曲率辐射 (Curvature Radiation)}

\n{在极强磁场（如中子星表面 $B \sim 10^{12}\text{ G}$）环境中，电子的动力学行为与弱场环境显著不同。同步辐射的高效冷却机制会导致垂直于磁力线的动量分量迅速衰减，使得电子被“束缚”在磁力线上滑行。此时，磁力线本身的弯曲成为了产生辐射的主导因素。}

\subsubsection{强磁场下的电子动力学限制}

\begin{theorem}[同步辐射冷却时标]
在强磁场中，相对论电子因同步辐射损耗垂直动量的寿命 $t_{sy}$ 极短。对于洛伦兹因子为 $\gamma$、垂直磁场分量为 $B_{\perp}$ 的电子，其寿命为（CGS 单位制）：
\begin{align}
t_{sy} &\approx 10^{9} \gamma^{-1} B_{\perp}^{-2} \, (\text{s}) \label{eq:synch_lifetime}
\end{align}
\end{theorem}

\begin{example}[中子星环境下的冷却]
对于典型的中子星参数，取电子 $\gamma = 1000$，磁场 $B = 10^{12}\text{ G}$（假设投射角适中，即 $B_{\perp} \sim B$），代入 \eqref{eq:synch_lifetime} 可得：
\begin{align}
t_{sy} &\approx 10^{9} \times (10^3)^{-1} \times (10^{12})^{-2} = 10^{-18} \, \text{s}
\end{align}
\end{example}

\begin{derivationnote}[物理图像：从同步辐射到曲率辐射的过渡]
极短的 $t_{sy}$ 意味着电子在极短时间内就会辐射掉垂直于磁力线的动量 $p_{\perp}$，导致其基态运动轨迹几乎完全沿磁力线平行（$p_{\parallel} \gg p_{\perp}$）。此时，同步辐射机制失效。然而，若磁力线本身具有曲率半径 $R_c$，沿磁力线运动的相对论电子将获得法向加速度 $a_{\perp} \approx v^2/R_c \approx c^2/R_c$。这种由空间几何曲率引起的加速辐射被称为\textbf{曲率辐射}。
\end{derivationnote}

\subsubsection{曲率辐射的频谱特征}

类似于同步辐射，我们首先定义描述辐射谱的关键频率参数。

\begin{definition}[基频与临界频率]
设电子沿曲率半径为 $R_c$ 的轨道运动。
1. **基频** $\nu_0$ 定义为轨道回旋频率：
\begin{align}
\nu_0 &= \frac{c}{2\pi R_c} \label{eq:nu_0}
\end{align}
2. **临界频率（或峰值频率）** $\nu_c$ 考虑了相对论效应带来的频率提升：
\begin{align}
\nu_c &= \frac{3}{2}\gamma^3 \nu_0 = \frac{3}{2}\gamma^3 \frac{c}{2\pi R_c} \label{eq:nu_c}
\end{align}
\end{definition}

\begin{theorem}[单粒子曲率辐射谱功率]
能量为 $\gamma m_e c^2$ 的电子产生的曲率辐射，其单位频率间隔的辐射功率 $P_{\nu}$ 为：
\begin{align}
P_{\nu} &= \frac{2\pi\sqrt{3}e^2 \nu_0 \gamma}{c} \frac{\nu}{\nu_c} \int_{\nu/\nu_c}^{\infty} K_{5/3}(t) dt \label{eq:curvature_power}
\end{align}
其中 $K_{5/3}(t)$ 为第二类修正贝塞尔函数。
\end{theorem}

\begin{proof}
（注：此公式结构与同步辐射完全一致，区别仅在于 $\nu_c$ 的物理定义不同。推导涉及推迟势的傅里叶变换及最速下降法近似，此处直接引用讲义给出的标准结果。）
\end{proof}

\subsubsection{低频渐近行为与同步辐射的对比}

\n{尽管曲率辐射与同步辐射的公式形式相似（均涉及 $K_{5/3}$ 积分），但在低频端（$\nu \ll \nu_c$），两者的物理依赖性存在本质区别。特别是曲率辐射的低频强度与电子能量 $\gamma$ 无关，这是识别辐射机制的重要特征。}

\begin{theorem}[低频谱指数与能量无关性]
在低频极限下（$\nu/\nu_c \ll 1$），曲率辐射谱呈现 $P_{\nu} \propto \nu^{1/3}$ 的幂律特征，且该辐射强度与电子洛伦兹因子 $\gamma$ 无关。
\end{theorem}

\begin{proof}
利用修正贝塞尔函数的小宗量渐近展开，当 $x \to 0$ 时，$\int_x^{\infty} K_{5/3}(t) dt \approx C x^{-2/3}$（其中 $C$ 为常数）。
将此代入 \eqref{eq:curvature_power}：
\begin{align}
P_{\nu} (\nu \ll \nu_c) &\propto \gamma \nu_0 \left( \frac{\nu}{\nu_c} \right) \left( \frac{\nu}{\nu_c} \right)^{-2/3} \nonumber \\
&= \gamma \nu_0 \left( \frac{\nu}{\nu_c} \right)^{1/3} \label{eq:low_freq_step1}
\end{align}
将 $\nu_c \propto \gamma^3 \nu_0$（见公式 \eqref{eq:nu_c}）代入上式：
\begin{align}
P_{\nu} &\propto \gamma \nu_0 \nu^{1/3} (\gamma^3 \nu_0)^{-1/3} \nonumber \\
&= \gamma \nu_0 \nu^{1/3} \gamma^{-1} \nu_0^{-1/3} \nonumber \\
&= \gamma^{0} \nu_0^{2/3} \nu^{1/3} \label{eq:gamma_independence}
\end{align}
可见，$\gamma$ 的幂次完全抵消。
\end{proof}

\begin{derivationnote}[与同步辐射的对比]
作为对比，对于同步辐射，其临界频率 $\nu_{c, \mathrm{syn}} \propto \gamma^2 B$。若重复上述推导：
\begin{align}
P_{\nu, \mathrm{syn}} &\propto \gamma (\nu / \gamma^2)^{1/3} \propto \gamma^{1/3} \nu^{1/3}
\end{align}
因此，\textbf{同步辐射的低频部分依然微弱地依赖于电子能量（$\propto \gamma^{1/3}$），而曲率辐射的低频部分完全独立于电子能量。}这一差异源于两种机制中临界频率对 $\gamma$ 的依赖律不同（$\gamma^3$ vs $\gamma^2$）。
\end{derivationnote}

\subsubsection{辐射的角分布}

\begin{theorem}[相对论聚束效应]
由于电子的超相对论运动（$\gamma \gg 1$），曲率辐射并非各向同性，而是高度集中在电子瞬时速度方向的狭窄锥体内。其半张角 $\theta$ 满足：
\begin{align}
\theta \sim \frac{1}{\gamma} \label{eq:beaming}
\end{align}
\end{theorem}

\begin{proof}
这是洛伦兹变换的直接推论。在电子静止系中发射的偶极辐射，经洛伦兹变换到实验室系后，光子动量主要集中在前进方向 $1/\gamma$ 的锥角内。
\end{proof}